\documentclass[cnatzke_thesis_proposal.tex]{subfiles}
\begin{document}

\chapter{The GRIFFIN Spectrometer at TRIUMF}

%------------------------------------------
\subsection{TRIUMF Rare Isotope Beam Facility}
%------------------------------------------
\begin{figure}[H]
  \begin{center}
    \includegraphics[width=0.95\linewidth]{triumf_cyclotron_construction_1972.jpg}
  \end{center}
  \caption{TRIUMF cyclotron during initial construction (1972) - Figure courtesy of TRIUMF.}
  \label{fig:triumf_cyclotron_1972}
\end{figure}

The proposed experiment will be performed at Canada's national laboratory for nuclear and particle physics, TRIUMF. 
TRIUMF is located in Vancouver, British Columbia, Canada and was established in 1968 by three universities, Simon Fraser University, the University of British Columbia (UBC), and the University of Victoria, as an experimental facility to meet needs the three universities could not individually provide. 
Since then the science program has grown to include nuclear physics, particle physics, molecular and material science, and nuclear medicine while providing research infrastructure and tools too large or complex for a single university to build, operate, or maintain. 
TRIUMF houses the world's largest cyclotron \cite{dilling_isac_2014} capable of accelerating protons up to 520 MeV before they are delivered to the experimental halls, via one of four beam lines, for direct use or to create secondary radioactive isotope beams of pions, muons, or radioactive isotopes. 
This work uses the GRIFFIN spectrometer housed in the ISAC-I experimental hall located on the TRIUMF-ISAC-I experimental hall. 

\begin{center}
\begin{figure}[H]
  \begin{center}
    \includegraphics[scale=.18]{isac_schematic.png}
  \end{center}
  \caption{Schematic view of the TRIUMF-ISAC facility \cite{dilling_isac_2014}, including the target ion-source, the high resolution mass separator, and various detection stations.}
  \label{fig:ISAC_HALL}
\end{figure}
\end{center}

%------------------------------------------
\subsection{The GRIFFIN Spectrometer}
%------------------------------------------

Large scale detector arrays for $\gamma$-ray measurements provide a powerful and robust tool for studying unstable nuclei through radioactive decay and nuclear spectroscopy at radioactive ion beam facilities \cite{garnsworthy_griffin_2019}. 
The Gamma-Ray Infrastructure For Fundamental Investigations of Nuclei (GRIFFIN) is a Compton-suppressed, high-efficiency $\gamma$-ray spectrometer composed of sixteen large-volume clover-type High Purity Germanium (HPGe) detectors located in the ISAC-I experimental hall at TRIUMF. 
Each clover contains four n-type HPGe crystals of 60 mm diameter before shaping and 90 mm in length which are tapered towards the front face to facilitate closer packing within the clover housing \cite{rizwan_characteristics_2016}.
The four crystals operating voltage is between 3.5 to 4 kV at 95 K and share the same cryostat cooled by a dewar of liquid nitrogen that is filled every 8 hours. 
The clover detectors, and any ancillary detector, are placed in an aluminum rhombicuboctahedron support structure where up to sixteen of the eighteen faces of the rhombicuboctahedron are covered by the HPGe clovers.
The remaining two faces are used for the incoming low-energy RIB from ISAC-I and the tape collection system that removes long-lived activity from the implantation chamber at the end of a measurement cycle \cite{garnsworthy_griffin_2019}. 
The western hemisphere of GRIFFIN is shown in Figure \ref{fig:griffin_westhemi} where the upstream beam line can be seen on the rightmost edge of the image. 

\begin{center}
  \begin{figure}[H]
    \begin{center}
      \includegraphics[width=0.75\textwidth]{griffin_clover_schematic.jpg}
    \end{center}
    \caption{Rendering of a GRIFFIN HPGe clover with the exterior dimensional tolerances of the aluminum crystal housing indicated. Figure from Ref.~\cite{rizwan_characteristics_2016}.}
    \label{fig:griffin_clover_schematic}
  \end{figure}
\end{center}

\begin{center}
  \begin{figure}[H]
    \begin{center}
      \includegraphics[width=0.75\textwidth]{griffin_westhemi.jpg}
    \end{center}
    \caption{Western hemisphere of GRIFFIN in optimized-peak-to-total mode taken during TRIUMF photowalk 2018. Image courtesy of TRIUMF.}
    \label{fig:griffin_westhemi}
  \end{figure}
\end{center}

Each GRIFFIN clover is surrounded by a set of high-efficiency bismuth germanate (BGO) active Compton suppression shields.
The BGO shields are made of a high-density, large atomic number ($Z$ = 83 for Bi) scintillator coupled to a photomultiplier tube (PMT) to create a detector ideally suited to measure escaping $\gamma$-rays due to its high $\gamma$-ray interaction cross-section. 
The Compton shields detect Compton-scattered $\gamma$-rays that have scattered out of a HPGe crystal, 511 keV $\gamma$-rays created via pair production, and provide passive shielding from background radiation originating outside the array. 
Singles from the suppression shields are used as a veto to reject incomplete scatter events in the HPGe crystals thereby increasing the sensitivity and signal-to-noise ratio of the array. 
The shields are composed of three individual components; a front plate, a side shield, and a back catcher, all of which can be moved independently providing freedom in the physical arrangement of the detector array. 

The GRIFFIN spectrometer has two distinct operating configurations based on the relative position of the BGO suppression shields and HPGe clover faces to the centre of the array. 
The first mode named "high-efficiency mode" is when the clovers are 110 mm radially distant from the centre of the array and the front BGO shields are pulled back, shown in the top panel of Figure~\ref{fig:griffin_modes}, to facilitate a large solid angle coverage and greater absolute detection efficiency. 
The second mode, called "optimized peak-to-total" mode shown in the bottom panel of Figure~\ref{fig:griffin_modes}, is when the HPGe clovers are 145 mm from the centre of the array and the front shield plates of the BGO suppressors are rolled forward to form a complete suppression shield around each HPGe clover. 

\begin{center}
  \begin{figure}[H]
    \begin{center}
      \includegraphics[width=0.95\textwidth]{griffin_modes.jpg}
    \end{center}
    \caption{Two possible configurations of the GRIFFIN spectrometer. The beam is delivered from the left. The Pb wall shielding the tape box can be seen on the right. Upper panel showing PACES in the upstream (left) half of the chamber and the fast scintillator in the downstream (right) chamber. In the ‘High-efficiency’ mode, the HPGe detectors are at 11 cm and the LaBr(Ce) detectors at 12.5 cm from the implantation point on the tape. Lower panel with SCEPTAR in the upstream and downstream chambers. In the ‘Optimized peak-to-total’ mode, the HPGe and LaBr(Ce) detectors are retracted to 14.5 cm and 13.5 cm, respectively, in order that they can be fully Compton and background suppressed with BGO shields. The 20 mm Delrin absorber is also installed in the lower panel. Figure taken from Ref. \cite{garnsworthy_griffin_2019}.}
    \label{fig:griffin_modes}
  \end{figure}
\end{center}

%------------------------------------------
\subsection{Data Acquisition}
\label{sec:data_acquisition}
%------------------------------------------
The GRIFFIN data acquisition system (DAQ) uses custom-designed digital electronic modules in a three tiered system \cite{garnsworthy_griffin_2017}. 
The modules were designed to handle high-rate data collection with each crystal counting at a rate of up to 50 kHz and process signals from the HPGe crystals, BGO suppression shields, and any ancillary detector used in conjunction with the clovers. 
The DAQ operates using three different levels of electronics; analogue-to-digital converters, secondary collection modules, and a primary collector module. The lowest level of the DAQ that processes raw detector outputs is made of GRIF-16 analogue-to-digital converters (ADC) that sample the raw analogue detector output at 100 MHz and convert the signal into a digital signal.
The ADCs pass the converted signal to the GRIF-C secondary collection modules that act as multiplexer for multiple input GRIF-16s. 
Finally, the output signal of the secondary collectors are passed to a single GRIF-C primary collector that filters the accumulated data from the lower levels according to user specified logic to accept or reject the singles to be written out. 
All three levels of the DAQ are synchronized to a GRIF-Clk Clock module that houses a Symmetricom Model SA.45s Chip-Sized Atomic Clock (CSAC) that produces a 10 MHz reference signal. 
The reference signal is used to generate a 50 MHz primary clock single that is fanned-out through secondary GRIF-Clk modules to all the collectors and digitizers in the system \cite{garnsworthy_griffin_2017}.

Events that pass the primary collector are written to disk using the Maximum Integrated Data Acquisition System (MIDAS) frontend \cite{ritt_midas_1997}, developed by TRIUMF and the Paul Scherrer Institute. 
Similar to the rest of the GRIFFIN DAQ MIDAS allows for high-rate data to be written to disk, and the data are written to disk using the GRIFFIN event format \cite{griffin_griffin_nodate}. 

%------------------------------------------
\subsection{Data Sorting}
%------------------------------------------
GRIFFIN data is written to disk according to the GRIFFIN event file format, a form of the general MIDAS file format (see Section~\ref{sec:data_acquisition}).
To properly analyze the data files they must first be unpacked using the analysis package GRSISort~\cite{bildstein_griffincollaborationgrsisort_2019}, a framework built on top of ROOT~\cite{brun_root-projectroot_2019}.
GRSISort converts the MIDAS files to ROOT files via a two stage process.
During the first unpacking stage GRSISort produces a fragment tree, a ROOT file containing time-ordered hits, then by applying an event-building algorithm GRSISort stores hits occurring within a user-specified time coincidence together in another ROOT file called an analysis tree.
In this work a 1 $\mu$s build window, the time-coincidence condition mentioned previously, was used.
The analysis trees can then be sorted into histograms via GRSISort or a custom sorting program compiled against the GRSISort libraries.

%------------------------------------------
\subsection{Standard GRIFFIN Calibrations and Data Conditioning}
%------------------------------------------
Before any investigation into two-photon decay detection can be performed the data has to be prepared and conditioned in a such a way as to maximize GRIFFIN's potential of measuring a two-photon transition. 
The techniques are processes discussed hereafter have been applied across as large range of GRIFFIN datasets and therefore are well understood. 
References~\cite{smith_gamma-gamma_2019}, \cite{garcia_absence_2020}, \cite{maclean_spectroscopy_2021}, and \cite{porzio_configuration_2021} all contain a more in-depth discussion of the common GRIFFIN data conditioning techniques.

%------------------------------------------
\subsubsection{GRIFFIN Energy Calibrations}
%------------------------------------------
To properly take advantage of GRIFFIN's high efficiency and large solid angle coverage data from the HPGe crystals are summed together throughout the analysis.
This summation of data requires a careful energy calibration to preserve high peak resolution of the detector array when summing data together from different crystals. 
The standard method for energy calibration uses known source data to calibrate the energy spectra for the crystals making up the sixteen clovers used in this experiment. 

The standard calibration technique uses transitions from well-studied calibrations sources: $^{60}$Co, $^{152}$Eu, $^{133}$Ba, and $^{56}$Co; where the peak energies are fit using a quadratic equation of the form: 

\begin{equation}
  E = a_0 + a_1 P + a_2 P^2
\end{equation}

where $a_i$ are the fitted coefficients, $P$ is the ADC pulse height, and $E$ is the energy of the $gamma$-ray. 
This quadratic calibration is applied on a crystal-by-crystal basis for all HPGe crystal present in the array. 
The effect of this calibration is shown in~\ref{fig:calibration_effect}. 

\begin{figure}[htbp]
  \centering
  \includegraphics[width=0.95\textwidth]{calibration_effect.png}
  \caption{Effect of calibration a GRIFFIN subrun. In this particular dataset clover 13, corresponding to HPGe crystals 49-52, was not present. Figure from~\cite{maclean_spectroscopy_2021}.
    \textit{Top:} 344 keV $\gamma$-ray emitted during the $\beta$ decay of $^{152}$Eu before the energy calibration is applied. 
    \textit{Bottom:} Calibrated 344 keV $\gamma$-ray emitted during the $\beta$ decay of $^{152}$Eu. 
  }
  \label{fig:calibration_effect}
\end{figure}

The $^{90}$Sr source data described in Section~\ref{sed:proof_of_concept_90sr} uses this same calibration technique applied on a run-by-run basis to optimize the energy resolution across the entire data-collection.

%------------------------------------------
\subsubsection{Angular Correlations}
%------------------------------------------
When no external field is present, excited nuclear states populated by decay processes generally have magnetic sub-states of equal population probability. 
Although the relaxation of each sub-state via $\gamma$-ray emission creates an anisotropic spatial distribution, the summation of all components from the equally populated magnetic sub-states creates and isotopic $\gamma$-ray spatial distribution when viewed from the lab frame, schematically shown in~\ref{fig:angular_correlation_substate}.

\begin{figure}[htbp]
  \centering
  \includegraphics[width=0.95\textwidth]{angular_correlation_state_isotropic.pdf}
  \caption{Anisotropic spatial distribution of depopulating $\gamma$-rays from individual magnetic substates. When summed over equally populated substates the distribution becomes isotropic in the lab frame. Figure adapted from \cite{maclean_spectroscopy_2021}.}
  \label{fig:angular_correlation_substate}
\end{figure}

If a nucleus emits multiple $\gamma$-rays sequentially, called a $\gamma$-cascade, the first $\gamma$-ray can generate an effective polarization of nucleus by defining a $z$-axis and therefore define the orientation of the magnetic sub-states. 
This orientation of states creates an unequal distribution of magnetic sub-states and an anisotropic spatial distribution of the depopulating $\gamma$-rays relative to the first $\gamma$-ray emitted in the cascade. 
This anisotropy is described by the function $W(\theta)$: 

\begin{equation}
  W(\theta) = \sum_{k, even} a_k P_k(\cos(\theta))
\end{equation}

where the coefficients, $a_k$, and the Legendre polynomials, $P_k$, are dependent on the spins and multipolarities for the $\gamma$ cascade of interest~\cite{krane_introductory_1987}. 
The sum only includes even $k$ values, to conserve parity, up to the maximum value $k_{max}$ equal to twice the lowest of the $\gamma$-ray multipolarities and the spin of the intermediate state involved in the cascade. 
The coefficients $a_k$ do not depend on the parity of the states and therefore angular correlation measurements cannot be used assign the parities of the states without other spectroscopic information. 
\ref{fig:angular_correlation_theory} shows the different angular correlations for a  $J_I^+ \rightarrow 2^+ \rightarrow 0^+$ with mixing ratio $\delta = 0$.

\begin{figure}[htbp]
  \centering
  \includegraphics[width=0.95\textwidth]{theoretical_ggac.png}
  \caption{Theoretical $\gamma-\gamma$ angular correlation for a $J_I^+ \rightarrow 2^+ \rightarrow 0^+$ cascade. All correlations shown have a mixing ratio of $\delta = 0$. Image from~\cite{maclean_spectroscopy_2021}.}
  \label{fig:angular_correlation_theory}
\end{figure}

Experimentally $\gamma-\gamma$ angular correlations provide a powerful tool for constraining and establishing spin assignments and multipole mixing ratios for transitions with mixed electromagnetic character. 

Proper $\gamma-\gamma$ angular correlations rely on using a multi-detector array to measure the intensity of $\gamma-\gamma$ emission at different relative angles, a task GRIFFIN is well suited for with its 52 distinct opening angles~\cite{smith_gamma-gamma_2019}.
The techniques used in analyzing $\gamma-\gamma$ angular correlations with GRIFFIN are well understood and described in~\cite{smith_gamma-gamma_2019}, ~\cite{smith_spectroscopic_2020}, ~\cite{maclean_high-precision_2020}, and~\cite{porzio_configuration_2021}.

%------------------------------------------
\subsubsection{Compton Suppression}
%------------------------------------------
In 2019 GRIFFIN received a full set of BGO Compton suppression shields increasing the peak-to-total ratio of collected data via Compton suppression.
In this works a suppression window, the maximum time difference between a HPGe and BGO hit, of 300 ns was used for the first two tests and 30ns was used for the third test. \ref{fig:bgo_gamma_timing} shows a 300 ns time window superimposed over a time-difference spectrum between a GRIFFIN HPGe crystal and a BGO shield. 
In all data collections any BGO shield surrounding the clover has been used as a veto to suppress signals from and of the four HPGe crystals that compose the clover. 
The spectrum shown in~\ref{fig:suppression_comparison_eu152} shows the effect of BGO-suppression on $^{152}$Eu source data along with the reject hits.

\begin{figure}[htbp]
  \centering
  \includegraphics[width=0.95\textwidth]{bgo_gamma_timing.png}
  \caption{Time difference spectrum for coincident $\gamma$-rays, one recorded in an HPGe crystal and the other in a BGO shield, in the same GRIFFIN event using a $^{152}$Eu source. The shaded red area indicates the time window used to perform the Compton suppression. Figure from~\cite{porzio_configuration_2021}.}
  \label{fig:bgo_gamma_timing}
\end{figure}

\begin{figure}[htbp]
  \centering
  \includegraphics[width=0.95\textwidth]{suppression_comparison_eu152.png}
  \caption{Comparison of a $^{152}$Eu source in all singles (black), suppressed singles (red), and events rejected by Compton suppression (red). Energy values of the most intense peaks are plotted in keV. Figure from~\cite{porzio_configuration_2021}.}
  \label{fig:suppression_comparison_eu152}
\end{figure}

%------------------------------------------
\subsection{Room Background}
\label{sec:room_background}
%------------------------------------------
The measurement of nuclear two-photon decay requires the registration of very rare decay phenomena and thus requires a good understanding of background radiation sources. 
One such background, spurious signals that can mimic the signals of interest, is natural $\gamma$-rays emitted primarily by the decay of three radiative isotopes: potassium-40 ($^{40}$K), uranium-238 ($^{238}$U), and thorium-208 ($^{208}$Th)~\cite[]{aksoy_elemental_1994}.
$^{40}$K decays directly to the stable argon-40 ($^{40}$Ar) emitting a 1460 keV gamma ray, and both $^{238}$U and $^{208}$Th decay through long decay chains eventually terminating in stable lead isotopes. 
Of all the $\gamma$-rays emitted in those long decay chains the most prominent are the 1765 keV and 2615 keV $\gamma$-rays emitted by bismuth-214 ($^{214}$Bi) and thalium-208 ($^{208}$Tl) respectively.

\begin{figure}[htbp]
  \centering
  \includegraphics[width=0.95\textwidth]{background_spectrum.png}
  \caption{$\gamma$-ray energy spectrum from 300 hours of GRIFFIN room background data. The major background peaks have been labeled.}
  \label{fig:background_spectrum}
\end{figure}

%------------------------------------------
\subsection{Sum Spectrum}
\label{sec:sum_spectrum}
%------------------------------------------

The two photons released during a two-photon decay event sum to the total energy of the transition, meaning the full energy of the transition should show up as a peak in a sum spectrum. 
A sum spectrum is generated by summing the energies of two incident photons together and recording their resultant sum in a histogram.
This is similar to the GRIFFIN addback algorithm, where the energy of all events detected within a single clover are summed together, but does not restrict the sum to singles events within the same clover; but rather, sums events together that were registered in different crystals regardless of if they were in the same clover of not. 
\ref{fig:singles_vs_sum_spectrum} shows the difference between a singles $\gamma$-ray energy spectrum and a sum spectrum.

\begin{figure}[htbp]
  \centering
  \includegraphics[width=0.95\textwidth]{singles_vs_sum_spectrum.png}
  \caption{Comparison of singles (black) and sum (green) spectra for 506 hours of $^{90}$Sr source collection.}
  \label{fig:singles_vs_sum_spectrum}
\end{figure}


% ------------------------------------
\end{document}
